Table~\ref{tab:combi:activity} shows activities for geometric objects:

\orbitertableofcommands{combinatorial_object_activity}

Some commands require extra options. 
Table~\ref{tab:combi:reporting} shows options for the report which is generated after a classification.
\orbitertableofcommands{objects_report_options}
For more details on classification, 
please see Section~\ref{sec:canonical:form:overview}.


\bigskip



The Edge curve is given by the equation
$$
X^4 -Y^4 -Z^4 +2f^2Y^2Z^2 +4fX^2YZ =0
$$
where $f$ is a primitive element of $\bbF_q.$
Let us pick $q=17.$
The next example creates the 
Edge curve in $\PG(2,17)$ using the primitive element $f=3$.
For more details, see Section~\ref{sec:algebraic:varieties}.
%The equation is encoded as a vector of coefficients, 
%using the ordering of monomials 
%as in Table~\ref{tab:PG2monomials}.

%\sourcecode{EDGE_CURVE_Q17_EQUATION}

%The following command computes the points 
%on the variety and writes them to a file. 
%The filename is \verb'Edge_q17.csv'.

%\sourcecode{EDGE_CURVE_Q17_AS_POINTS}


%\sourcecode{FILE_Q17}


\sourcecode{Edge_curve_17}


The following command computes the line type of the Edge curve:

\sourcecode{Edge_curve_17_line_type}

The line type is 
$$
( 4^6, 2^{30}, 1^{132}, 0^{139} )
$$
This means that there are 6 4-secants, 
30 2-secants, 132 tangent lines, and 139 external lines to the curve.






